\documentclass[12pt,a4paper]{article}
\usepackage[UTF8]{ctex}
\usepackage{amsmath,amssymb}
\usepackage{geometry}
\usepackage{listings}
\usepackage{xcolor}
\usepackage{tcolorbox}
\usepackage{enumitem}
\usepackage{hyperref}
\usepackage{float}

\geometry{margin=2.5cm}

% 定义颜色
\definecolor{lightgray}{RGB}{240,240,240}
\definecolor{myblue}{RGB}{70,130,180}
\definecolor{mygreen}{RGB}{34,139,34}
\definecolor{myred}{RGB}{220,20,60}
\definecolor{myorange}{RGB}{255,140,0}
\definecolor{myyellow}{RGB}{218,165,32}
\definecolor{mypurple}{RGB}{128,0,128}

% C++ 代码高亮设置
\lstset{
	language=C++,
	basicstyle=\ttfamily\small,
	keywordstyle=\color{myblue}\bfseries,
	commentstyle=\color{gray},
	stringstyle=\color{teal},
	numbers=left,
	numberstyle=\tiny\color{gray},
	stepnumber=1,
	numbersep=5pt,
	backgroundcolor=\color{lightgray!30},
	frame=single,
	breaklines=true,
	breakatwhitespace=false,
	tabsize=4,
	showspaces=false,
	showstringspaces=false,
	showtabs=false,
	captionpos=b,
	morekeywords={long long, vector, pair, unordered\_map, unordered\_set, queue, priority\_queue, stack}
}

% 自定义盒子环境
\newtcolorbox{bluebox}{
	colback=myblue!10!white,
	colframe=myblue!70!black,
	fonttitle=\bfseries,
	arc=4pt
}

\newtcolorbox{greenbox}{
	colback=mygreen!10!white,
	colframe=mygreen!70!black,
	fonttitle=\bfseries,
	arc=4pt
}

\newtcolorbox{redbox}{
	colback=myred!10!white,
	colframe=myred!70!black,
	fonttitle=\bfseries,
	arc=4pt
}

\newtcolorbox{orangebox}{
	colback=myorange!10!white,
	colframe=myorange!70!black,
	fonttitle=\bfseries,
	arc=4pt
}

\newtcolorbox{yellowbox}{
	colback=myyellow!10!white,
	colframe=myyellow!70!black,
	fonttitle=\bfseries,
	arc=4pt
}

\newtcolorbox{purplebox}{
	colback=mypurple!10!white,
	colframe=mypurple!70!black,
	fonttitle=\bfseries,
	arc=4pt
}

\title{\Huge \textbf{BFS与DFS搜索算法完全指南}}
\author{算法学习笔记}
\date{\today}

\begin{document}
	
	\maketitle
	
	\tableofcontents
	\newpage
	
	\section{搜索算法基础概念}
	
	\subsection{什么是搜索算法？}
	
	搜索算法是计算机科学中最基础、最重要的算法之一。它用于在数据集合中查找特定信息或遍历整个数据结构。
	
	\begin{bluebox}
		\textbf{搜索算法的本质}
		
		搜索算法的核心是\textbf{状态空间探索}：
		\begin{itemize}
			\item 每个问题可以看作一个\textbf{状态图}
			\item 图中的节点代表\textbf{状态}
			\item 图中的边代表\textbf{状态转移}
			\item 搜索算法就是在这个图中找到从初始状态到目标状态的路径
		\end{itemize}
	\end{bluebox}
	
	\subsection{两种核心搜索策略}
	
	\begin{yellowbox}
		\textbf{DFS vs BFS 形象对比}
		
		\begin{itemize}
			\item \textbf{DFS（深度优先搜索）}：像走迷宫，一条路走到黑，走不通再回头
			\begin{itemize}
				\item 比喻：一个人探索迷宫，一直往前走直到无路可走，然后返回上一个岔路口
			\end{itemize}
			\item \textbf{BFS（广度优先搜索）}：像水波扩散，一层一层向外扩展
			\begin{itemize}
				\item 比喻：一群人同时从起点出发，每一步所有的人都向外走一步
			\end{itemize}
		\end{itemize}
	\end{yellowbox}
	
	\subsection{时间复杂度与空间复杂度对比}
	
	\begin{center}
		\begin{tabular}{|c|c|c|c|}
			\hline
			\textbf{算法} & \textbf{时间复杂度} & \textbf{空间复杂度} & \textbf{数据结构} \\
			\hline
			DFS & $O(V + E)$ & $O(h)$ 递归栈深度 & 栈(Stack) \\
			\hline
			BFS & $O(V + E)$ & $O(w)$ 最大层宽度 & 队列(Queue) \\
			\hline
		\end{tabular}
	\end{center}
	
	其中 $V$ 是顶点数，$E$ 是边数，$h$ 是最大深度，$w$ 是最大宽度。
	
	\section{深度优先搜索（DFS）}
	
	\subsection{DFS的核心思想}
	
	DFS遵循"深入优先"的原则，尽可能深地搜索树或图的分支。
	
	\begin{greenbox}
		\textbf{DFS的三个关键步骤}
		
		\begin{enumerate}
			\item \textbf{选择}：从当前节点选择一个未访问的邻接节点
			\item \textbf{深入}：移动到该节点，继续深入探索
			\item \textbf{回溯}：当无路可走时，返回上一个节点，尝试其他分支
		\end{enumerate}
	\end{greenbox}
	
	\subsection{DFS的两种实现方式}
	
	\subsubsection{递归实现（最常用）}
	
	递归实现DFS是最自然的方式，利用系统栈来保存状态。
	
	\begin{lstlisting}[caption={DFS递归实现 - 图的遍历}]
		#include <iostream>
		#include <vector>
		using namespace std;
		
		const int MAXN = 100005;
		vector<int> graph[MAXN];  // 邻接表存图
		bool visited[MAXN];       // 标记节点是否访问过
		
		// DFS递归实现
		void dfs(int u) {
			// 1. 处理当前节点
			visited[u] = true;
			cout << "访问节点: " << u << endl;
			
			// 2. 遍历所有邻接节点
			for (int v : graph[u]) {
				// 3. 如果未访问，深入探索
				if (!visited[v]) {
					dfs(v);  // 递归深入
				}
			}
			// 4. 回溯（自动发生，递归返回即回溯）
		}
		
		// 使用示例
		int main() {
			int n, m;  // n个节点，m条边
			cin >> n >> m;
			
			// 建图
			for (int i = 0; i < m; i++) {
				int u, v;
				cin >> u >> v;
				graph[u].push_back(v);
				graph[v].push_back(u);  // 无向图
			}
			
			// 从节点1开始DFS
			dfs(1);
			
			return 0;
		}
	\end{lstlisting}
	
	\subsubsection{迭代实现（使用栈）}
	
	使用显式栈来模拟递归过程，避免递归深度过大导致栈溢出。
	
	\begin{lstlisting}[caption={DFS迭代实现 - 使用栈}]
		#include <iostream>
		#include <vector>
		#include <stack>
		using namespace std;
		
		const int MAXN = 100005;
		vector<int> graph[MAXN];
		bool visited[MAXN];
		
		// DFS迭代实现
		void dfs_iterative(int start) {
			stack<int> stk;
			stk.push(start);
			
			while (!stk.empty()) {
				int u = stk.top();
				stk.pop();
				
				// 可能在栈中有重复节点，需要再次检查
				if (visited[u]) continue;
				
				visited[u] = true;
				cout << "访问节点: " << u << endl;
				
				// 将所有未访问的邻接节点入栈
				for (int v : graph[u]) {
					if (!visited[v]) {
						stk.push(v);
					}
				}
			}
		}
		
		int main() {
			int n = 5;
			graph[1] = {2, 3};
			graph[2] = {1, 4, 5};
			graph[3] = {1};
			graph[4] = {2};
			graph[5] = {2};
			
			dfs_iterative(1);
			
			return 0;
		}
	\end{lstlisting}
	
	\subsection{DFS的应用场景}
	
	\begin{bluebox}
		\textbf{DFS的典型应用}
		
		\begin{enumerate}
			\item \textbf{连通性问题}：判断图的连通分量数量
			\item \textbf{路径搜索}：查找从起点到终点的所有路径
			\item \textbf{拓扑排序}：有向无环图的线性排序
			\item \textbf{环检测}：判断图中是否存在环
			\item \textbf{回溯算法}：排列、组合、子集等问题
			\item \textbf{树的遍历}：前序、中序、后序遍历
		\end{enumerate}
	\end{bluebox}
	
	\subsection{DFS经典问题：全排列}
	
	\begin{lstlisting}[caption={DFS回溯 - 全排列问题}]
		#include <iostream>
		#include <vector>
		using namespace std;
		
		vector<int> path;           // 当前路径（当前排列）
		vector<bool> used;          // 标记数字是否使用过
		vector<vector<int>> result; // 存储所有结果
		
		// 生成1~n的所有排列
		void backtrack(int n) {
			// 终止条件：路径长度等于n，得到一个完整排列
			if (path.size() == n) {
				result.push_back(path);
				// 打印排列
				for (int num : path) {
					cout << num << " ";
				}
				cout << endl;
				return;
			}
			
			// 尝试每个数字
			for (int i = 1; i <= n; i++) {
				if (!used[i]) {
					// 选择
					used[i] = true;
					path.push_back(i);
					
					// 深入
					backtrack(n);
					
					// 回溯（撤销选择）
					path.pop_back();
					used[i] = false;
				}
			}
		}
		
		int main() {
			int n = 3;
			used.resize(n + 1, false);
			
			cout << "1~3的所有排列:" << endl;
			backtrack(n);
			
			return 0;
		}
		/* 输出:
		1 2 3
		1 3 2
		2 1 3
		2 3 1
		3 1 2
		3 2 1
		*/
	\end{lstlisting}
	
	\section{广度优先搜索（BFS）}
	
	\subsection{BFS的核心思想}
	
	BFS遵循"广度优先"的原则，逐层访问节点。
	
	\begin{greenbox}
		\textbf{BFS的关键特征}
		
		\begin{enumerate}
			\item \textbf{层次性}：按照距离起点的层数逐层访问
			\item \textbf{最短路径}：在无权图中能找到最短路径
			\item \textbf{队列实现}：使用队列保证先进先出的顺序
			\item \textbf{非递归}：天然适合迭代实现
		\end{enumerate}
	\end{greenbox}
	
	\subsection{BFS的标准模板}
	
	\begin{lstlisting}[caption={BFS标准模板 - 图的层序遍历}]
		#include <iostream>
		#include <vector>
		#include <queue>
		using namespace std;
		
		const int MAXN = 100005;
		vector<int> graph[MAXN];
		bool visited[MAXN];
		int dist[MAXN];  // 记录从起点到每个节点的距离
		
		// BFS实现
		void bfs(int start) {
			queue<int> q;
			
			// 1. 将起点入队
			q.push(start);
			visited[start] = true;
			dist[start] = 0;
			
			// 2. 开始BFS
			while (!q.empty()) {
				int u = q.front();
				q.pop();
				
				cout << "访问节点: " << u 
				<< ", 距离: " << dist[u] << endl;
				
				// 3. 遍历所有邻接节点
				for (int v : graph[u]) {
					if (!visited[v]) {
						// 4. 未访问的节点入队
						visited[v] = true;
						dist[v] = dist[u] + 1;  // 更新距离
						q.push(v);
					}
				}
			}
		}
		
		int main() {
			int n = 5;
			graph[1] = {2, 3};
			graph[2] = {1, 4, 5};
			graph[3] = {1};
			graph[4] = {2};
			graph[5] = {2};
			
			cout << "BFS从节点1开始:" << endl;
			bfs(1);
			
			return 0;
		}
		/* 输出:
		访问节点: 1, 距离: 0
		访问节点: 2, 距离: 1
		访问节点: 3, 距离: 1
		访问节点: 4, 距离: 2
		访问节点: 5, 距离: 2
		*/
	\end{lstlisting}
	
	\subsection{BFS的应用场景}
	
	\begin{bluebox}
		\textbf{BFS的典型应用}
		
		\begin{enumerate}
			\item \textbf{最短路径}：无权图中求最短路径
			\item \textbf{层次遍历}：树的层序遍历
			\item \textbf{连通分量}：统计连通块数量
			\item \textbf{迷宫问题}：找最短逃生路径
			\item \textbf{单词接龙}：词语变换的最短序列
			\item \textbf{社交网络}：计算好友关系的度数
		\end{enumerate}
	\end{bluebox}
	
	\subsection{BFS经典问题：迷宫最短路径}
	
	\begin{lstlisting}[caption={BFS - 迷宫最短路径}]
		#include <iostream>
		#include <vector>
		#include <queue>
		using namespace std;
		
		struct Node {
			int x, y;      // 坐标
			int steps;     // 从起点到这里的步数
		};
		
		// 四个方向：上、下、左、右
		int dx[] = {-1, 1, 0, 0};
		int dy[] = {0, 0, -1, 1};
		
		int shortestPath(vector<vector<int>>& maze, 
		int startX, int startY, 
		int endX, int endY) {
			int n = maze.size();
			int m = maze[0].size();
			
			vector<vector<bool>> visited(n, vector<bool>(m, false));
			queue<Node> q;
			
			// 起点入队
			q.push({startX, startY, 0});
			visited[startX][startY] = true;
			
			while (!q.empty()) {
				Node cur = q.front();
				q.pop();
				
				// 到达终点
				if (cur.x == endX && cur.y == endY) {
					return cur.steps;
				}
				
				// 尝试四个方向
				for (int i = 0; i < 4; i++) {
					int nx = cur.x + dx[i];
					int ny = cur.y + dy[i];
					
					// 检查边界和障碍
					if (nx >= 0 && nx < n && ny >= 0 && ny < m 
					&& maze[nx][ny] == 0 && !visited[nx][ny]) {
						visited[nx][ny] = true;
						q.push({nx, ny, cur.steps + 1});
					}
				}
			}
			
			return -1;  // 无法到达
		}
		
		int main() {
			// 0表示通路，1表示障碍
			vector<vector<int>> maze = {
				{0, 0, 0, 0},
				{0, 1, 0, 0},
				{0, 1, 0, 1},
				{0, 0, 0, 0}
			};
			
			int steps = shortestPath(maze, 0, 0, 3, 3);
			cout << "最短路径步数: " << steps << endl;
			// 输出: 6
			
			return 0;
		}
	\end{lstlisting}
	
	\section{DFS与BFS的选择策略}
	
	\begin{orangebox}
		\textbf{如何选择DFS还是BFS？}
		
		\begin{center}
			\begin{tabular}{|c|c|c|}
				\hline
				\textbf{问题特征} & \textbf{推荐算法} & \textbf{原因} \\
				\hline
				求最短路径/最少步数 & BFS & 逐层扩展，首次到达即最短 \\
				\hline
				求所有可能解/路径 & DFS & 回溯机制天然适合枚举 \\
				\hline
				树/图很深 & BFS & 避免递归栈溢出 \\
				\hline
				树/图很宽 & DFS & BFS队列会占用大量空间 \\
				\hline
				检测连通性 & 都可以 & 性能相近 \\
				\hline
				拓扑排序 & DFS & 容易实现 \\
				\hline
				状态空间很大 & DFS+剪枝 & 内存占用小 \\
				\hline
			\end{tabular}
		\end{center}
	\end{orangebox}
	
	\section{经典例题详解}
	
	\subsection{例题1：岛屿数量（基础DFS/BFS）}
	
	\textbf{题目描述：}给定一个由 '1'（陆地）和 '0'（水）组成的二维网格，计算网格中岛屿的数量。岛屿总是被水包围，并且每座岛屿只能由水平方向或竖直方向上相邻的陆地连接形成。
	
	\begin{lstlisting}[caption={例题1：岛屿数量 - DFS解法}]
		#include <iostream>
		#include <vector>
		using namespace std;
		
		// DFS淹没岛屿
		void dfs(vector<vector<char>>& grid, int i, int j) {
			int n = grid.size();
			int m = grid[0].size();
			
			// 边界检查
			if (i < 0 || i >= n || j < 0 || j >= m || grid[i][j] == '0') {
				return;
			}
			
			// 将当前陆地标记为水（淹没）
			grid[i][j] = '0';
			
			// 向四个方向扩展
			dfs(grid, i - 1, j);  // 上
			dfs(grid, i + 1, j);  // 下
			dfs(grid, i, j - 1);  // 左
			dfs(grid, i, j + 1);  // 右
		}
		
		int numIslands(vector<vector<char>>& grid) {
			if (grid.empty()) return 0;
			
			int n = grid.size();
			int m = grid[0].size();
			int count = 0;
			
			for (int i = 0; i < n; i++) {
				for (int j = 0; j < m; j++) {
					if (grid[i][j] == '1') {
						count++;  // 发现新岛屿
						dfs(grid, i, j);  // 淹没整个岛屿
					}
				}
			}
			
			return count;
		}
		
		int main() {
			vector<vector<char>> grid = {
				{'1', '1', '0', '0', '0'},
				{'1', '1', '0', '0', '0'},
				{'0', '0', '1', '0', '0'},
				{'0', '0', '0', '1', '1'}
			};
			
			cout << "岛屿数量: " << numIslands(grid) << endl;
			// 输出: 3
			
			return 0;
		}
	\end{lstlisting}
	
	\subsection{例题2：岛屿数量 - BFS解法}
	
	\begin{lstlisting}[caption={例题1变体：岛屿数量 - BFS解法}]
		#include <iostream>
		#include <vector>
		#include <queue>
		using namespace std;
		
		int numIslandsBFS(vector<vector<char>>& grid) {
			if (grid.empty()) return 0;
			
			int n = grid.size();
			int m = grid[0].size();
			int count = 0;
			
			int dx[] = {-1, 1, 0, 0};
			int dy[] = {0, 0, -1, 1};
			
			for (int i = 0; i < n; i++) {
				for (int j = 0; j < m; j++) {
					if (grid[i][j] == '1') {
						count++;
						
						// BFS淹没岛屿
						queue<pair<int, int>> q;
						q.push({i, j});
						grid[i][j] = '0';
						
						while (!q.empty()) {
							auto [x, y] = q.front();
							q.pop();
							
							for (int k = 0; k < 4; k++) {
								int nx = x + dx[k];
								int ny = y + dy[k];
								
								if (nx >= 0 && nx < n && ny >= 0 && ny < m 
								&& grid[nx][ny] == '1') {
									grid[nx][ny] = '0';
									q.push({nx, ny});
								}
							}
						}
					}
				}
			}
			
			return count;
		}
	\end{lstlisting}
	
	\subsection{例题3：单词接龙（BFS求最短路径）}
	
	\textbf{题目描述：}给定两个单词（beginWord和endWord）和一个字典，找到从beginWord到endWord的最短转换序列的长度。每次转换只能改变一个字母，且转换后的单词必须在字典中。
	
	\begin{lstlisting}[caption={例题3：单词接龙 - BFS}]
		#include <iostream>
		#include <vector>
		#include <queue>
		#include <unordered_set>
		#include <unordered_map>
		using namespace std;
		
		int ladderLength(string beginWord, string endWord, 
		vector<string>& wordList) {
			// 将字典转换为集合，方便查找
			unordered_set<string> wordSet(wordList.begin(), wordList.end());
			
			// 如果终点不在字典中，无法转换
			if (wordSet.find(endWord) == wordSet.end()) {
				return 0;
			}
			
			// BFS
			queue<string> q;
			q.push(beginWord);
			
			int steps = 1;  // 包含起始单词
			
			while (!q.empty()) {
				int size = q.size();
				
				// 处理当前层的所有节点
				for (int i = 0; i < size; i++) {
					string word = q.front();
					q.pop();
					
					// 尝试改变每个位置的字母
					for (int j = 0; j < word.length(); j++) {
						char original = word[j];
						
						// 尝试26个字母
						for (char c = 'a'; c <= 'z'; c++) {
							word[j] = c;
							
							// 找到终点
							if (word == endWord) {
								return steps + 1;
							}
							
							// 如果在字典中，加入队列继续搜索
							if (wordSet.find(word) != wordSet.end()) {
								q.push(word);
								wordSet.erase(word);  // 删除防止重复访问
							}
						}
						
						// 恢复原字母
						word[j] = original;
					}
				}
				
				steps++;  // 层数增加
			}
			
			return 0;  // 无法转换
		}
		
		int main() {
			string beginWord = "hit";
			string endWord = "cog";
			vector<string> wordList = {"hot", "dot", "dog", "lot", "log", "cog"};
			
			int result = ladderLength(beginWord, endWord, wordList);
			cout << "最短转换序列长度: " << result << endl;
			// 输出: 5 (hit -> hot -> dot -> dog -> cog)
			
			return 0;
		}
	\end{lstlisting}
	
	\subsection{例题4：课程表（拓扑排序 - DFS）}
	
	\textbf{题目描述：}现在你总共有n门课需要选，记为0到n-1。在选修某些课程之前需要先修其他课程。给定课程总量和先决条件，判断是否可能完成所有课程的学习。
	
	\begin{lstlisting}[caption={例题4：课程表 - DFS检测环}]
		#include <iostream>
		#include <vector>
		using namespace std;
		
		// 课程表问题 - DFS检测环
		bool hasCycle(int course, vector<vector<int>>& graph, 
		vector<int>& visited) {
			// 1: 正在访问（在当前DFS路径中）
			// 2: 已经访问完成
			if (visited[course] == 1) return true;   // 发现环
			if (visited[course] == 2) return false;  // 已经处理过
			
			visited[course] = 1;  // 标记为正在访问
			
			// 访问所有后继课程
			for (int next : graph[course]) {
				if (hasCycle(next, graph, visited)) {
					return true;
				}
			}
			
			visited[course] = 2;  // 标记为访问完成
			return false;
		}
		
		bool canFinish(int numCourses, vector<vector<int>>& prerequisites) {
			// 构建图
			vector<vector<int>> graph(numCourses);
			for (auto& pre : prerequisites) {
				graph[pre[1]].push_back(pre[0]);
			}
			
			// 访问状态：0=未访问, 1=访问中, 2=已访问
			vector<int> visited(numCourses, 0);
			
			// 对每个节点进行DFS
			for (int i = 0; i < numCourses; i++) {
				if (visited[i] == 0) {
					if (hasCycle(i, graph, visited)) {
						return false;  // 有环，无法完成
					}
				}
			}
			
			return true;  // 无环，可以完成
		}
		
		int main() {
			int n = 4;
			vector<vector<int>> prerequisites = {{1, 0}, {2, 0}, {3, 1}, {3, 2}};
			
			cout << "能否完成所有课程: " 
			<< (canFinish(n, prerequisites) ? "是" : "否") << endl;
			// 输出: 是
			
			return 0;
		}
	\end{lstlisting}
	
	\subsection{例题5：课程表 - BFS拓扑排序}
	
	\begin{lstlisting}[caption={例题4变体：课程表 - BFS拓扑排序}]
		#include <iostream>
		#include <vector>
		#include <queue>
		using namespace std;
		
		bool canFinishBFS(int numCourses, vector<vector<int>>& prerequisites) {
			// 构建图和入度数组
			vector<vector<int>> graph(numCourses);
			vector<int> inDegree(numCourses, 0);
			
			for (auto& pre : prerequisites) {
				graph[pre[1]].push_back(pre[0]);
				inDegree[pre[0]]++;
			}
			
			// 将所有入度为0的节点入队
			queue<int> q;
			for (int i = 0; i < numCourses; i++) {
				if (inDegree[i] == 0) {
					q.push(i);
				}
			}
			
			int count = 0;  // 可以完成的课程数
			
			// BFS拓扑排序
			while (!q.empty()) {
				int course = q.front();
				q.pop();
				count++;
				
				// 减少所有后继课程的入度
				for (int next : graph[course]) {
					inDegree[next]--;
					if (inDegree[next] == 0) {
						q.push(next);
					}
				}
			}
			
			return count == numCourses;  // 所有课程都完成了
		}
	\end{lstlisting}
	
	\section{进阶题目详解}
	
	\subsection{进阶题1：被围绕的区域（DFS/BFS）}
	
	\textbf{题目描述：}给定一个二维矩阵，包含'X'和'O'。找到所有被'X'围绕的区域，并将这些区域里所有的'O'用'X'填充。边界上的'O'不会被填充。
	
	\textbf{解题思路：}
	\begin{itemize}
		\item 核心思想：反过来思考，找到不被包围的'O'（边界上的'O'及其连通的'O'）
		\item 从边界开始DFS/BFS，标记所有与边界连通的'O'
		\item 未被标记的'O'就是被包围的，改为'X'
	\end{itemize}
	
	\begin{lstlisting}[caption={进阶题1：被围绕的区域}]
		#include <iostream>
		#include <vector>
		using namespace std;
		
		void dfs(vector<vector<char>>& board, int i, int j) {
			int n = board.size();
			int m = board[0].size();
			
			if (i < 0 || i >= n || j < 0 || j >= m || board[i][j] != 'O') {
				return;
			}
			
			// 标记为临时字符，表示不被包围
			board[i][j] = '#';
			
			dfs(board, i - 1, j);
			dfs(board, i + 1, j);
			dfs(board, i, j - 1);
			dfs(board, i, j + 1);
		}
		
		void solve(vector<vector<char>>& board) {
			if (board.empty()) return;
			
			int n = board.size();
			int m = board[0].size();
			
			// 处理四条边界
			for (int i = 0; i < n; i++) {
				dfs(board, i, 0);      // 左边界
				dfs(board, i, m - 1);  // 右边界
			}
			
			for (int j = 0; j < m; j++) {
				dfs(board, 0, j);      // 上边界
				dfs(board, n - 1, j);  // 下边界
			}
			
			// 转换
			for (int i = 0; i < n; i++) {
				for (int j = 0; j < m; j++) {
					if (board[i][j] == 'O') {
						board[i][j] = 'X';  // 被包围的改为X
					} else if (board[i][j] == '#') {
						board[i][j] = 'O';  // 不被包围的恢复
					}
				}
			}
		}
		
		int main() {
			vector<vector<char>> board = {
				{'X', 'X', 'X', 'X'},
				{'X', 'O', 'O', 'X'},
				{'X', 'X', 'O', 'X'},
				{'X', 'O', 'X', 'X'}
			};
			
			solve(board);
			
			cout << "处理后的矩阵:" << endl;
			for (auto& row : board) {
				for (char c : row) {
					cout << c << " ";
				}
				cout << endl;
			}
			
			return 0;
		}
	\end{lstlisting}
	
	\subsection{进阶题2：扫雷游戏（BFS + 模拟）}
	
	\textbf{题目描述：}给定一个代表游戏面板的二维字符矩阵，点击一个位置，根据扫雷规则揭示该位置。
	
	\textbf{解题思路：}
	\begin{itemize}
		\item 如果点击到地雷（M），直接改为X并结束
		\item 如果点击到空格（E），用BFS/DFS展开：
		\begin{itemize}
			\item 计算周围地雷数量
			\item 如果有地雷，显示数字
			\item 如果没地雷，显示B，并递归展开周围8个格子
		\end{itemize}
	\end{itemize}
	
	\begin{lstlisting}[caption={进阶题2：扫雷游戏}]
		#include <iostream>
		#include <vector>
		#include <queue>
		using namespace std;
		
		vector<vector<char>> updateBoard(vector<vector<char>>& board, 
		vector<int>& click) {
			int n = board.size();
			int m = board[0].size();
			int x = click[0], y = click[1];
			
			// 点击到地雷
			if (board[x][y] == 'M') {
				board[x][y] = 'X';
				return board;
			}
			
			// 8个方向
			int dx[] = {-1, -1, -1, 0, 0, 1, 1, 1};
			int dy[] = {-1, 0, 1, -1, 1, -1, 0, 1};
			
			// BFS
			queue<pair<int, int>> q;
			q.push({x, y});
			
			// 标记已处理，避免重复入队
			vector<vector<bool>> inQueue(n, vector<bool>(m, false));
			inQueue[x][y] = true;
			
			while (!q.empty()) {
				auto [cx, cy] = q.front();
				q.pop();
				
				// 计算周围地雷数量
				int mines = 0;
				vector<pair<int, int>> neighbors;
				
				for (int i = 0; i < 8; i++) {
					int nx = cx + dx[i];
					int ny = cy + dy[i];
					
					if (nx >= 0 && nx < n && ny >= 0 && ny < m) {
						if (board[nx][ny] == 'M') {
							mines++;
						} else if (!inQueue[nx][ny]) {
							neighbors.push_back({nx, ny});
						}
					}
				}
				
				if (mines > 0) {
					// 有地雷，显示数字
					board[cx][cy] = '0' + mines;
				} else {
					// 没有地雷，递归展开
					board[cx][cy] = 'B';
					for (auto [nx, ny] : neighbors) {
						q.push({nx, ny});
						inQueue[nx][ny] = true;
					}
				}
			}
			
			return board;
		}
		
		int main() {
			vector<vector<char>> board = {
				{'E', 'E', 'E', 'E', 'E'},
				{'E', 'E', 'M', 'E', 'E'},
				{'E', 'E', 'E', 'E', 'E'},
				{'E', 'E', 'E', 'E', 'E'}
			};
			vector<int> click = {3, 0};
			
			auto result = updateBoard(board, click);
			
			cout << "扫雷结果:" << endl;
			for (auto& row : result) {
				for (char c : row) {
					cout << c << " ";
				}
				cout << endl;
			}
			
			return 0;
		}
	\end{lstlisting}
	
	\subsection{进阶题3：太平洋大西洋水流问题}
	
	\textbf{题目描述：}给定一个m x n的非负整数矩阵表示地形高度。水可以从高处向低处流。求所有同时能流向太平洋和大西洋的坐标。太平洋在左边界和上边界，大西洋在右边界和下边界。
	
	\begin{lstlisting}[caption={进阶题3：太平洋大西洋水流问题}]
		#include <iostream>
		#include <vector>
		using namespace std;
		
		void dfs(vector<vector<int>>& heights, vector<vector<bool>>& ocean,
		int i, int j, int prevHeight) {
			int n = heights.size();
			int m = heights[0].size();
			
			// 检查边界和高度条件
			if (i < 0 || i >= n || j < 0 || j >= m 
			|| ocean[i][j] || heights[i][j] < prevHeight) {
				return;
			}
			
			// 标记为可达
			ocean[i][j] = true;
			
			// 向四个方向流动（注意：水从低往高流是反过来的）
			dfs(heights, ocean, i - 1, j, heights[i][j]);
			dfs(heights, ocean, i + 1, j, heights[i][j]);
			dfs(heights, ocean, i, j - 1, heights[i][j]);
			dfs(heights, ocean, i, j + 1, heights[i][j]);
		}
		
		vector<vector<int>> pacificAtlantic(vector<vector<int>>& heights) {
			if (heights.empty()) return {};
			
			int n = heights.size();
			int m = heights[0].size();
			
			vector<vector<bool>> pacific(n, vector<bool>(m, false));
			vector<vector<bool>> atlantic(n, vector<bool>(m, false));
			
			// 从太平洋边界开始DFS（左边界和上边界）
			for (int i = 0; i < n; i++) {
				dfs(heights, pacific, i, 0, INT_MIN);
			}
			for (int j = 0; j < m; j++) {
				dfs(heights, pacific, 0, j, INT_MIN);
			}
			
			// 从大西洋边界开始DFS（右边界和下边界）
			for (int i = 0; i < n; i++) {
				dfs(heights, atlantic, i, m - 1, INT_MIN);
			}
			for (int j = 0; j < m; j++) {
				dfs(heights, atlantic, n - 1, j, INT_MIN);
			}
			
			// 收集同时能流向两大洋的点
			vector<vector<int>> result;
			for (int i = 0; i < n; i++) {
				for (int j = 0; j < m; j++) {
					if (pacific[i][j] && atlantic[i][j]) {
						result.push_back({i, j});
					}
				}
			}
			
			return result;
		}
		
		int main() {
			vector<vector<int>> heights = {
				{1, 2, 2, 3, 5},
				{3, 2, 3, 4, 4},
				{2, 4, 5, 3, 1},
				{6, 7, 1, 4, 5},
				{5, 1, 1, 2, 4}
			};
			
			auto result = pacificAtlantic(heights);
			
			cout << "同时能流向太平洋和大西洋的坐标:" << endl;
			for (auto& pos : result) {
				cout << "(" << pos[0] << "," << pos[1] << ") ";
			}
			cout << endl;
			
			return 0;
		}
	\end{lstlisting}
	
	\subsection{进阶题4：最短路径访问所有节点（状态压缩BFS）}
	
	\textbf{题目描述：}存在一个由n个节点组成的无向连通图，节点编号从0到n-1。找出最短的路径，使得从任意节点出发，访问完所有节点。可以以任意节点作为起点，任意节点作为终点，可以多次访问节点和重复经过边。
	
	\begin{lstlisting}[caption={进阶题4：最短路径访问所有节点 - BFS+状态压缩}]
		#include <iostream>
		#include <vector>
		#include <queue>
		#include <tuple>
		using namespace std;
		
		int shortestPathLength(vector<vector<int>>& graph) {
			int n = graph.size();
			
			// BFS队列：(当前节点, 访问状态, 路径长度)
			queue<tuple<int, int, int>> q;
			
			// visited[node][state] 防止重复访问
			vector<vector<bool>> visited(n, vector<bool>(1 << n, false));
			
			// 可以从任意节点开始
			for (int i = 0; i < n; i++) {
				int state = 1 << i;
				q.push({i, state, 0});
				visited[i][state] = true;
			}
			
			int target = (1 << n) - 1;  // 所有节点都访问过的状态
			
			while (!q.empty()) {
				auto [node, state, dist] = q.front();
				q.pop();
				
				// 所有节点都访问了
				if (state == target) {
					return dist;
				}
				
				// 访问相邻节点
				for (int next : graph[node]) {
					int newState = state | (1 << next);
					
					if (!visited[next][newState]) {
						visited[next][newState] = true;
						q.push({next, newState, dist + 1});
					}
				}
			}
			
			return -1;  // 不应该到达这里
		}
		
		int main() {
			vector<vector<int>> graph = {
				{1, 2, 3},
				{0},
				{0},
				{0}
			};
			
			cout << "最短路径长度: " << shortestPathLength(graph) << endl;
			// 输出: 4
			
			graph = {{1}, {0, 2, 4}, {1, 3, 4}, {2}, {1, 2}};
			cout << "最短路径长度: " << shortestPathLength(graph) << endl;
			// 输出: 4
			
			return 0;
		}
	\end{lstlisting}
	
	\subsection{进阶题5：01矩阵（多源BFS）}
	
	\textbf{题目描述：}给定一个由0和1组成的矩阵，找出每个元素到最近的0的距离。
	
	\begin{lstlisting}[caption={进阶题5：01矩阵 - 多源BFS}]
		#include <iostream>
		#include <vector>
		#include <queue>
		using namespace std;
		
		vector<vector<int>> updateMatrix(vector<vector<int>>& mat) {
			int n = mat.size();
			int m = mat[0].size();
			
			vector<vector<int>> dist(n, vector<int>(m, INT_MAX));
			queue<pair<int, int>> q;
			
			// 将所有0入队（多源BFS）
			for (int i = 0; i < n; i++) {
				for (int j = 0; j < m; j++) {
					if (mat[i][j] == 0) {
						dist[i][j] = 0;
						q.push({i, j});
					}
				}
			}
			
			int dx[] = {-1, 1, 0, 0};
			int dy[] = {0, 0, -1, 1};
			
			// BFS
			while (!q.empty()) {
				auto [x, y] = q.front();
				q.pop();
				
				for (int k = 0; k < 4; k++) {
					int nx = x + dx[k];
					int ny = y + dy[k];
					
					if (nx >= 0 && nx < n && ny >= 0 && ny < m 
					&& dist[nx][ny] > dist[x][y] + 1) {
						dist[nx][ny] = dist[x][y] + 1;
						q.push({nx, ny});
					}
				}
			}
			
			return dist;
		}
		
		int main() {
			vector<vector<int>> mat = {
				{0, 0, 0},
				{0, 1, 0},
				{1, 1, 1}
			};
			
			auto result = updateMatrix(mat);
			
			cout << "到最近0的距离:" << endl;
			for (auto& row : result) {
				for (int d : row) {
					cout << d << " ";
				}
				cout << endl;
			}
			// 输出:
			// 0 0 0
			// 0 1 0
			// 1 2 1
			
			return 0;
		}
	\end{lstlisting}
	
	\section{搜索算法的优化技巧}
	
	\subsection{剪枝优化}
	
	\begin{bluebox}
		\textbf{DFS剪枝策略}
		
		\begin{enumerate}
			\item \textbf{可行性剪枝}：当前状态不可能达到目标时提前返回
			\item \textbf{最优性剪枝}：当前状态不如已知最优解时提前返回
			\item \textbf{对称性剪枝}：利用问题的对称性减少搜索空间
			\item \textbf{记忆化搜索}：避免重复搜索相同状态
		\end{enumerate}
	\end{bluebox}
	
	\begin{lstlisting}[caption={DFS剪枝示例：N皇后问题}]
		#include <iostream>
		#include <vector>
		using namespace std;
		
		vector<vector<string>> result;
		
		bool isValid(vector<string>& board, int row, int col, int n) {
			// 检查列
			for (int i = 0; i < row; i++) {
				if (board[i][col] == 'Q') return false;
			}
			
			// 检查左上方对角线
			for (int i = row - 1, j = col - 1; i >= 0 && j >= 0; i--, j--) {
				if (board[i][j] == 'Q') return false;
			}
			
			// 检查右上方对角线
			for (int i = row - 1, j = col + 1; i >= 0 && j < n; i--, j++) {
				if (board[i][j] == 'Q') return false;
			}
			
			return true;
		}
		
		void backtrack(vector<string>& board, int row, int n) {
			if (row == n) {
				result.push_back(board);
				return;
			}
			
			for (int col = 0; col < n; col++) {
				// 可行性剪枝：只尝试合法位置
				if (isValid(board, row, col, n)) {
					board[row][col] = 'Q';
					backtrack(board, row + 1, n);
					board[row][col] = '.';
				}
			}
		}
		
		vector<vector<string>> solveNQueens(int n) {
			vector<string> board(n, string(n, '.'));
			backtrack(board, 0, n);
			return result;
		}
		
		int main() {
			int n = 4;
			auto solutions = solveNQueens(n);
			
			cout << n << "皇后问题的解法数量: " << solutions.size() << endl;
			for (auto& solution : solutions) {
				for (auto& row : solution) {
					cout << row << endl;
				}
				cout << endl;
			}
			
			return 0;
		}
	\end{lstlisting}
	
	\subsection{双向BFS}
	
	\begin{lstlisting}[caption={双向BFS示例：单词接龙优化}]
		#include <iostream>
		#include <vector>
		#include <unordered_set>
		#include <unordered_map>
		using namespace std;
		
		int ladderLengthBiBFS(string beginWord, string endWord, 
		vector<string>& wordList) {
			unordered_set<string> dict(wordList.begin(), wordList.end());
			if (!dict.count(endWord)) return 0;
			
			unordered_set<string> beginSet{beginWord};
			unordered_set<string> endSet{endWord};
			int steps = 1;
			
			while (!beginSet.empty() && !endSet.empty()) {
				// 总是扩展较小的集合
				if (beginSet.size() > endSet.size()) {
					swap(beginSet, endSet);
				}
				
				unordered_set<string> nextSet;
				
				for (string word : beginSet) {
					for (int i = 0; i < word.length(); i++) {
						char original = word[i];
						
						for (char c = 'a'; c <= 'z'; c++) {
							word[i] = c;
							
							if (endSet.count(word)) {
								return steps + 1;
							}
							
							if (dict.count(word)) {
								nextSet.insert(word);
								dict.erase(word);
							}
						}
						
						word[i] = original;
					}
				}
				
				beginSet = nextSet;
				steps++;
			}
			
			return 0;
		}
	\end{lstlisting}
	
	\subsection{记忆化搜索（Memoization）}
	
	\begin{lstlisting}[caption={记忆化搜索示例：最长递增路径}]
		#include <iostream>
		#include <vector>
		#include <algorithm>
		using namespace std;
		
		int dfs(vector<vector<int>>& matrix, vector<vector<int>>& memo, 
		int i, int j) {
			if (memo[i][j] != 0) {
				return memo[i][j];  // 记忆化：直接返回已计算的结果
			}
			
			int n = matrix.size();
			int m = matrix[0].size();
			int maxLen = 1;
			
			int dx[] = {-1, 1, 0, 0};
			int dy[] = {0, 0, -1, 1};
			
			for (int k = 0; k < 4; k++) {
				int ni = i + dx[k];
				int nj = j + dy[k];
				
				if (ni >= 0 && ni < n && nj >= 0 && nj < m 
				&& matrix[ni][nj] > matrix[i][j]) {
					maxLen = max(maxLen, 1 + dfs(matrix, memo, ni, nj));
				}
			}
			
			memo[i][j] = maxLen;  // 存储结果
			return maxLen;
		}
		
		int longestIncreasingPath(vector<vector<int>>& matrix) {
			if (matrix.empty()) return 0;
			
			int n = matrix.size(), m = matrix[0].size();
			vector<vector<int>> memo(n, vector<int>(m, 0));
			int result = 0;
			
			for (int i = 0; i < n; i++) {
				for (int j = 0; j < m; j++) {
					result = max(result, dfs(matrix, memo, i, j));
				}
			}
			
			return result;
		}
		
		int main() {
			vector<vector<int>> matrix = {
				{9, 9, 4},
				{6, 6, 8},
				{2, 1, 1}
			};
			
			cout << "最长递增路径长度: " << longestIncreasingPath(matrix) << endl;
			// 输出: 4 (1 -> 2 -> 6 -> 9)
			
			return 0;
		}
	\end{lstlisting}
	
	\section{总结与对比}
	
	\begin{greenbox}
		\textbf{BFS vs DFS 核心对比总结}
		
		\begin{center}
			\begin{tabular}{|c|c|c|}
				\hline
				\textbf{特性} & \textbf{BFS} & \textbf{DFS} \\
				\hline
				数据结构 & 队列 & 栈（或递归） \\
				\hline
				遍历顺序 & 层次遍历 & 深度遍历 \\
				\hline
				最短路径 & 能找到最优解 & 不一定 \\
				\hline
				空间复杂度 & $O(w)$ - 宽度 & $O(h)$ - 深度 \\
				\hline
				实现难度 & 较简单 & 递归简单，迭代较难 \\
				\hline
				适用场景 & 最短路、层序 & 枚举所有解、回溯 \\
				\hline
				栈溢出风险 & 无（迭代） & 有（递归深度大） \\
				\hline
				剪枝难度 & 较难 & 容易 \\
				\hline
			\end{tabular}
		\end{center}
	\end{greenbox}
	
	\begin{purplebox}
		\textbf{面试高频考点}
		
		\begin{enumerate}
			\item \textbf{岛屿系列}：岛屿数量、最大岛屿面积、岛屿周长
			\item \textbf{迷宫系列}：最短路径、墙与门、逃离迷宫
			\item \textbf{单词系列}：单词接龙、单词搜索
			\item \textbf{排列组合}：全排列、子集、组合总和
			\item \textbf{图遍历}：克隆图、课程表、网络延迟
			\item \textbf{矩阵搜索}：被围绕区域、太平洋大西洋水流
			\item \textbf{状态搜索}：打开转盘锁、滑动谜题
		\end{enumerate}
	\end{purplebox}
	
	\section{练习题推荐}
	
	\begin{enumerate}[leftmargin=*]
		\item \textbf{DFS基础}：
		\begin{itemize}
			\item 岛屿的最大面积（LeetCode 695）
			\item 省份数量（LeetCode 547）
			\item 组合总和（LeetCode 39）
		\end{itemize}
		\item \textbf{BFS基础}：
		\begin{itemize}
			\item 二叉树的层序遍历（LeetCode 102）
			\item 墙与门（LeetCode 286）
			\item 腐烂的橘子（LeetCode 994）
		\end{itemize}
		\item \textbf{进阶挑战}：
		\begin{itemize}
			\item 打开转盘锁（LeetCode 752）
			\item 滑动谜题（LeetCode 773）
			\item 公交路线（LeetCode 815）
		\end{itemize}
		\item \textbf{困难模式}：
		\begin{itemize}
			\item 最短超级串（LeetCode 943）
			\item 逃离大迷宫（LeetCode 1036）
			\item 访问所有节点的最短路径（LeetCode 847）
		\end{itemize}
	\end{enumerate}
	
\end{document}